\documentclass{report}
\usepackage{amsmath,amssymb}
\setlength{\parindent}{0mm}
\setlength{\parskip}{1em}
\begin{document}
\begin{center}
\rule{6in}{1pt} \
{\large Christopher M. Siefert \\
{\bf Won't you be my neighbor? A look at neighbor discovery for algebraic multigrid}}

Sandia National Laboratories \\ P O Box 5800 \\ MS 1322 \\ Albuquerque \\ NM 87185-1322
\\
{\tt csiefer@sandia.gov}\end{center}

At the core of many MPI-based codes, including almost all MPI-based PDE codes,
is the idea of a ``neighbor'' --- the list processors to whom
point-to-point messages get sent and from whom they are received.
Thus, one of the very first bits of parallel code an implementer must
develop (or borrow from elsewhere) is code that does arbitrary
neighbor discovery. More specifically, this code answers the
question, ``Who owns a particular degree of freedom?'' This code is necessary for
things as mundane as assembling sparse matrices and performing
matrix-vector products. Thus a non-trivial
literature has developed on performing arbitrary neighbor discovery
with minimum storage and message cost. With such a tool in the
proverbial toolbox, it is not uncommon for the arbitrary neighbor
discovery code to be used again and again --- even for problems where
additional information is available to make neighbor discovery \textit{not} so
arbitrary.

In this talk we will focus neighbor discovery in two (related)
computational kernels involved in the setup phase of algebraic
multigrid --- migration of sparse matrices between processors and sparse
matrix-matrix multiplication. We will show that for these kernels,
passing additional information during point-to-point communication
can allow one to perform neighbor discovery more efficiently.
For these kernels, it can be performed using a number of
messages and amount of storage independent of the number of processors
(assuming that the number of neighbors is bounded independently).


\end{document}
